\section{Results}
\label{sec:results}

\subsection{Predictive quality}
\label{sec:quality}

Table~\ref{tab:main} is the headline. Recursion alone is lowest-loss. MoR beats
MoE by $0.4329$ nats and beats MoRE by $0.1216$. Both recursive arms beat MoE
decisively, and the composition lands between them---the premise that routing
would \emph{add} to recursion in domain is not supported.

\begin{table}[t]
\centering
\small
\caption{Canonical results, 5 seeds each, mean$\,\pm\,$std. Primary metric is
final-epoch validation task loss; \emph{best} is the minimum over evaluated
epochs and equals \emph{last} on all 15 cells, so no model selection occurred.
$\Delta$ floor is the margin below the bigram baseline of 4.9849 nats
(positive \emph{beats} it). FLOPs are inference FLOPs/token at each arm's
measured mean depth, from \texttt{results/flops\_language.json}. Pairwise, on the
primary metric: MoRE$-$MoE $-0.3113$ ($d\,{-}15.9$, $p\,0.0079$); MoRE$-$MoR
$+0.1216$ ($d\,{+}4.7$, $p\,0.0079$); MoE$-$MoR $+0.4329$ ($d\,{+}14.7$,
$p\,0.0079$). Exact two-sided randomization test, floor $p_{\min} = 0.00794$;
every headline comparison sits on it.}
\label{tab:main}
\vspace{2pt}
\begin{tabular}{lrrrrrr}
\toprule
& Params & val task loss & PPL & $\Delta$ floor & FLOPs/tok & items/s \\
\midrule
MoE  & 5{,}584{,}908 & 3.9483\,$\pm$\,0.0242 & 51.9 & $+1.0366$ & 5.25\,M & 724 \\
MoR  & 5{,}581{,}063 & \textbf{3.5154\,$\pm$\,0.0340} & \textbf{33.6} & $+1.4695$ & 47.6\,M & 170 \\
MoRE & 5{,}584{,}908 & 3.6370\,$\pm$\,0.0134 & 38.0 & $+1.3480$ & 11.5\,M & 170 \\
\bottomrule
\end{tabular}
\end{table}

Two features of the table deserve to be read together. First, MoE and MoRE are
parameter-\emph{identical} at 5{,}584{,}908 while MoR sits $0.069\%$ lower at
5{,}581{,}063, the residual being the router MoR does not have---and that residual
runs \emph{against} the headline, since the arm with fewer parameters is the one
that wins. Second, MoRE has by far the tightest seed spread (std 0.0134 against
MoR's 0.0340), and we flag that we do not know what to make of it: it is the one
axis on which the composition is unambiguous, and the runs that exist cannot
separate a genuinely stabilising effect from the simpler account that a model whose
halting has collapsed to a constant has less to vary across seeds.

Fig.~\ref{fig:depth-curve} and the forced-depth sweep sharpen this. Both arms scale
with forced depth up to a point, in opposite directions:

\begin{itemize}[leftmargin=1.4em,itemsep=1.5pt,topsep=2pt]
\item \textbf{MoR} improves from 7.36 nats at depth 1 to a minimum of 4.06 at
  depth 4, then \emph{degrades} to 4.99 at depth 7. It has a genuine interior
  optimum.
\item \textbf{MoRE} improves monotonically all the way to depth 7 (3.70), never
  turning over within the budget. It wants more depth than it has.
\end{itemize}

\begin{figure}[t]
\centering
\includegraphics[width=0.92\linewidth]{fig_depth_curve.pdf}
\caption{Validation loss against forced uniform recursion depth, five seeds. Lines
are the forced-depth sweep; stars mark each arm's \emph{learned} policy at its own
measured mean depth. MoR's learned policy sits far below its own best uniform depth
(3.515 vs 4.056); MoRE's sits just below its forced-depth-7 endpoint (3.637 vs
3.697). The dotted line is MoE, which has no recursion. This is the paper's central
contrast: adaptive depth buys MoR a great deal and MoRE very little.}
\label{fig:depth-curve}
\end{figure}

\subsection{Is recursion actually adaptive?}
\label{sec:halting}

The mechanism claim rests on comparing each learned policy against a \emph{uniform}
policy at the same compute, which is the cheapest competitor that ignores its input
and therefore the minimum bar an adaptive-depth claim must clear. Two comparisons
matter, and they give different answers for the two arms.

\textbf{Against the best uniform depth anywhere} (the strongest competitor): MoR's
learned policy at 3.5154 beats its best forced depth (4.0558 at $d{=}4$) by
\textbf{0.540 nats}. MoRE's learned policy at 3.6370 beats its best forced depth
(3.6971 at $d{=}7$) by \textbf{0.060 nats}.

\textbf{Against a depth-matched uniform policy}---forced to exactly the arm's own
measured mean depth of 6.955, so the comparison isolates \emph{where} the compute
goes rather than \emph{how much} there is: MoR's adaptive policy beats forced-$d7$
(4.9863) by \textbf{1.471 nats}; MoRE's beats forced-$d7$ (3.6971) by
\textbf{0.060 nats}.

The depth-matched comparison is the one that licenses a mechanism claim, because
best-depth comparison conflates adaptivity with the average compute spent. On it,
both arms' halting is doing something real and measurable---but MoR's buys 1.47
nats where MoRE's buys 0.06, a factor of roughly 24. The honest summary is not that
MoRE's depth mechanism is inert; it is that adding routing left it almost entirely
without effect.

The exit statistics say why. MoR early-exits $18.8\% \pm 15.1\%$ of validation
tokens with mean depth $6.730$ and a halt remainder of $0.175$;
MoRE force-exits $97.3\% \pm 0.8\%$ at the ceiling with mean depth $6.955$ and
remainder $0.655$. MoRE's halting has collapsed to ``run to budget''.

\textbf{The depth--difficulty correlation is reported next to those rates because
it is only interpretable alongside them.} MoR's per-token depth correlates with
per-token model loss at $\rho = +0.2985 \pm 0.0654$, positive on \textbf{5/5}
seeds. MoRE's is $-0.0230 \pm 0.0115$---but with $97.3\%$ of tokens exiting at the
ceiling and a per-token depth standard deviation of 0.43, that correlation is
computed against a near-constant. We report it as \emph{uninformative}, not as
evidence that MoRE allocates depth badly. This distinction is the paper's most
load-bearing methodological point: a near-zero correlation with a constant is not
a finding about allocation.

MoR's early-exit rate is the seed-unstable quantity ($6.4\%$--$42.2\%$ across
seeds), and its instability is informative. The two seeds that exit most are the
two with the \emph{worst} loss, and the three that exit least are the best. So the
\emph{alignment} of depth with difficulty is robust while the \emph{amount} of
compute skipped is seed-unstable and costs loss---which is why the mean is reported
with its standard deviation rather than as a single depth budget.

\subsection{Where the composition does pay: robustness and efficiency}
\label{sec:ood}

Out of domain the ordering inverts. Table~\ref{tab:evals} evaluates all three arms
on the same checkpoints, and scores each against the inference FLOPs it costs.

\begin{table}[t]
\centering
\small
\caption{Evaluation battery, five seeds. Every cell is
\emph{score (difference vs MoE, inference FLOPs$\times$ vs MoE, difference per
FLOPs-multiplier)}; lower is better for loss rows, higher for accuracy rows.
``diff/$\times$'' is how much score each unit of extra compute bought, so higher is
more efficient. \textbf{These are exploratory}: they sit outside both declared
families, are uncorrected, and carry no significance verdict.}
\label{tab:evals}
\vspace{2pt}
\begin{tabular}{llll}
\toprule
Eval & MoE & MoR & MoRE \\
\midrule
Val loss (nats) & 3.948 (1.00$\times$) & 3.515 ($-0.433$, 9.07$\times$, 0.048) & 3.637 ($-0.311$, 2.19$\times$, \textbf{0.142}) \\
Test loss (nats) & 3.947 (1.00$\times$) & 3.513 ($-0.435$, 9.07$\times$, 0.048) & 3.621 ($-0.326$, 2.19$\times$, \textbf{0.149}) \\
OOD loss (nats) & 8.272 (1.00$\times$) & 7.802 ($-0.470$, 9.07$\times$, 0.052) & \textbf{7.573} ($-0.699$, 2.19$\times$, \textbf{0.319}) \\
BLiMP (acc) & 56.63 (1.00$\times$) & \textbf{58.74} ($+2.11$, 9.07$\times$, 0.233) & 58.36 ($+1.73$, 2.19$\times$, \textbf{0.791}) \\
EWoK (acc) & 49.62 (1.00$\times$)$^{*}$ & 50.33 ($+0.72$, 9.07$\times$, 0.079)$^{*}$ & 49.95 ($+0.33$, 2.19$\times$, 0.152)$^{*}$ \\
\bottomrule
\end{tabular}
\\[3pt]
{\footnotesize $^{*}$EWoK is within noise of the 50\% chance floor for every arm,
so its rankings are not meaningful---a scale limitation, reported for completeness.}
\end{table}

Three things follow.

\textbf{MoRE is the most robust arm out of domain.} On held-out Pile blocks
\citep{gao2020pile} it reaches 7.573 nats against MoR's 7.802 and MoE's 8.272---a
$0.699$-nat advantage over MoE against MoR's $0.470$. The in-domain ordering
inverts: the composition is the \emph{worst} recursive arm in domain and the
\emph{best} out of it. That is a real and reportable property, and it is the one
respect in which composing the two mechanisms does something neither does alone.

\textbf{It gets there far more cheaply.} MoR attains the lowest in-domain loss but
pays $9.07\times$ MoE's inference FLOPs; MoRE pays $2.19\times$. Per unit of extra
compute MoRE buys $0.142$ nats against MoR's $0.048$---roughly $3\times$ the
efficiency in its gain over MoE, and $6\times$ on the OOD row. FLOPs, not
wall-clock, is the compute measure; the throughput column is reported so the
latency cost stays visible, and it shows MoRE no faster than MoR in wall-clock
(170 items/s each) despite doing $4\times$ less arithmetic, because its per-token
expert dispatch is latency-bound.

\textbf{Grammatical competence is real but modest, and routing does not buy it.}
On BLiMP \citep{warstadt2020blimp}, 67{,}000 minimal pairs, both recursive arms beat
MoE (58.74\% and 58.36\% against 56.63\%) and MoR is not distinguishable from MoRE
($p = 0.71$). Every arm is above the 50\% floor, so recursion buys measurable
syntactic sensitivity at this scale; adding routing does not add more of it. The
gains concentrate on long-distance agreement phenomena---anaphor agreement and
determiner--noun agreement---which is where a shared iterative block would be
expected to help. On EWoK \citep{leong2023ewok} all three arms sit at chance, which
is a scale limitation and not a finding.

\subsection{Stratified loss: the advantage concentrates on hard tokens}
\label{sec:stratified}

Fig.~\ref{fig:stratified} groups validation loss by the frequency decile of the
\emph{target} token---difficulty being a property of what must be predicted, not of
the input---and by position within the block.

Both recursive arms improve most where prediction is hardest. The MoRE$-$MoE
advantage grows monotonically from $-0.147$ nats on the rarest-decile tokens to
$-0.374$ on the most frequent; MoR's from $-0.188$ to $-0.527$. Neither
architecture's gain is a uniform shift, which is what a pure capacity effect would
look like.

\begin{figure}[t]
\centering
\includegraphics[width=0.98\linewidth]{fig_stratified.pdf}
\caption{(a) Validation loss by target-token frequency decile. (b) Each recursive
arm's advantage over MoE, which widens as tokens get harder. Exploratory,
uncorrected.}
\label{fig:stratified}
\end{figure}

\textbf{A prediction this does not support.} If recursion were doing what the
adaptive-computation premise suggests---spending steps where context is thin---the
recursive arms should gain most at the \emph{start} of a block. They do the
opposite: MoRE's advantage over MoE is $-0.047$ nats in the first 8 positions and
$-0.359$ in the last 128; MoR's goes from $-0.250$ to $-0.465$. Both recursive arms
gain most where context is \emph{longest}. We report this because it disconfirms a
mechanism we would otherwise have asserted, and because it is consistent with the
capacity account in \S\ref{sec:mechanism}: a shared block iterated is accumulating
representation, not compensating for missing context.

\begin{figure}[t]
\centering
\includegraphics[width=0.98\linewidth]{fig_ood.pdf}
\caption{The ordering inverts out of domain: MoRE is worst in domain (a) and best
on the Pile (b). Same three entities and colours in both panels. Exploratory,
uncorrected.}
\label{fig:ood}
\end{figure}
